\chapter{A Probability Toolkit for Lawyers}
\label{app:toolkit}

This appendix collects the definitions and results used in the book, with short proofs or derivations where they are instructive. It is meant as a reference, not a course; readers who want a fuller treatment should consult \citet{lindley2006}, \citet{kadane2011} or \citet{aitken_taroni2004}.

\section{Probability}

\subsection{The three rules}

For propositions $A$, $B$ and background information $K$:
\begin{enumerate}[label=(\roman*)]
\item Convexity: $0 \le \Prob(A\given K) \le 1$, and $\Prob(A\given K) = 1$ if $K$ implies $A$.
\item Addition: if $A$ and $B$ are mutually exclusive given $K$, $\Prob(A \text{ or } B\given K) = \Prob(A\given K) + \Prob(B\given K)$.
\item Multiplication: $\Prob(A\text{ and } B\given K) = \Prob(A\given K)\,\Prob(B\given A, K)$.
\end{enumerate}
Everything else follows. The complement rule $\Prob(\bar A\given K) = 1 - \Prob(A\given K)$ follows from (i) and (ii). The law of total probability,
\[
\Prob(E\given K) = \sum_i \Prob(E\given H_i, K)\,\Prob(H_i\given K)
\]
for mutually exclusive and exhaustive $H_i$, follows from (ii) and (iii).

\subsection{Bayes' rule}

Applying the multiplication rule in both orders to $\Prob(H\text{ and }E\given K)$ gives
\[
\Prob(H\given E, K) = \frac{\Prob(E\given H, K)\,\Prob(H\given K)}{\Prob(E\given K)} .
\]
In odds form, for two hypotheses $\Hp$ and $\Hd$,
\[
\frac{\Prob(\Hp\given E,K)}{\Prob(\Hd\given E,K)} = \frac{\Prob(E\given\Hp,K)}{\Prob(E\given \Hd, K)}\times \frac{\Prob(\Hp\given K)}{\Prob(\Hd\given K)} .
\]
Posterior odds equal the likelihood ratio times the prior odds. When $\Hp$ and $\Hd$ are complements, odds $o$ convert to probability $o/(1+o)$.

\subsection{Independence}

$A$ and $B$ are independent given $K$ if $\Prob(A\given B, K) = \Prob(A\given K)$. Independence is always relative to background information: two witnesses' reports are dependent unconditionally, since both depend on what happened, but may be independent given what happened. One error in \emph{Collins} and in \emph{Sally Clark} was to multiply probabilities of features that are dependent.

\section{Evidence}

\subsection{Likelihood ratio and weight}

The likelihood ratio of evidence $E$ for $\Hp$ against $\Hd$ is $\LR = \Prob(E\given\Hp)/\Prob(E\given\Hd)$ (the background information $K$ is suppressed from here on). The weight of evidence in decibans is $10\log_{10}\LR$. Items that are independent given $\Hp$ and given $\Hd$ add their weights. Table~\ref{tab:decibans} gives conversions.

\begin{table}[ht]
\centering\small
\begin{tabular}{rrr}
\toprule
Likelihood ratio & Decibans & Posterior from prior $1/2$\\
\midrule
$1$ & $0$ & $0.500$\\
$2$ & $3$ & $0.667$\\
$10$ & $10$ & $0.909$\\
$100$ & $20$ & $0.990$\\
$1{,}000$ & $30$ & $0.999$\\
$10^{6}$ & $60$ & $0.999999$\\
\bottomrule
\end{tabular}
\caption{Likelihood ratios, weights in decibans, and posterior probabilities from even prior odds.}
\label{tab:decibans}
\end{table}

\subsection{Composite hypotheses}

When a hypothesis has unknown parameters $\theta$ with prior $\pi(\theta\given H)$, its marginal likelihood is $m(E\given H) = \int \Prob(E\given\theta,H)\pi(\theta\given H)\,d\theta$, and the Bayes factor is the ratio of marginal likelihoods (written $m_p(E)$ and $m_d(E)$ in Chapter~\ref{ch:bf}).

\subsection{The island formula}

If $N$ people, including the suspect $S$, could have left the stain, with prior probabilities $\pi_i$ of being the source and random match probabilities $\gamma_i$, and $S$ matches, then
\[
\Prob(S \text{ is the source}\given\text{match}) = \frac{\pi_S}{\pi_S + \sum_{i\ne S}\pi_i \gamma_i}.
\]
With uniform priors and a common $\gamma$ this is $1/(1 + (N-1)\gamma)$, as in Chapter~\ref{ch:odds}.

\subsection{Uncertain match probabilities}

If the random match probability $\gamma$ is uncertain with mean $\mu$ and variance $\sigma^2$, the effective match probability given one observed match is $\E(\gamma^2)/\E(\gamma) = \mu + \sigma^2/\mu$.

\subsection{Beta binomial prediction}

If a feature has frequency $\gamma \sim \text{Beta}(a,b)$ and has been seen $k$ times in $n$ items, the predictive probability that the next item has it is $(k+a)/(n+a+b)$.

\section{Decisions}

\subsection{Expected utility}

Given actions $d$, states $\theta$ with posterior $\Prob(\theta\given E)$, and utility $U(d,\theta)$, the optimal action maximises $\E[U(d,\theta)\given E] = \sum_\theta U(d,\theta)\Prob(\theta\given E)$. Equivalently it minimises expected loss.

\subsection{Thresholds}

With two actions and two states, and zero loss for correct decisions, act for $\Hp$ (convict, decision $c$) when $\Prob(\Hp\given E) > p^\ast = L(c,\bar G)/(L(c,\bar G) + L(a,G))$, where $G$ is $\Hp$, $\bar G$ is not $G$, $L(c,\bar G)$ is the loss from acting for $\Hp$ when $\Hd$ is true (in a criminal case, a false conviction) and $L(a,G)$ the loss from acquitting (decision $a$) when $\Hp$ is true. This is \eqref{eq:threshold}; Table~\ref{tab:thresholds} gives values.

\begin{table}[ht]
\centering\small
\begin{tabular}{lcc}
\toprule
Threshold $p^\ast$ & Loss ratio $L(c,\bar G)/L(a,G)$ & Threshold odds\\
\midrule
$0.50$ & $1$ & $1$\\
$0.75$ & $3$ & $3$\\
$0.90$ & $9$ & $9$\\
$0.95$ & $19$ & $19$\\
$0.99$ & $99$ & $99$\\
\bottomrule
\end{tabular}
\caption{Posterior thresholds and the loss ratios they imply. The loss ratio equals the threshold odds $p^\ast/(1-p^\ast)$, which is why the last two columns coincide.}
\label{tab:thresholds}
\end{table}

\subsection{Value of information}

The expected value of information from observing $X$ before deciding is the expected maximum expected utility after observing $X$ minus the maximum expected utility now. It is never negative, and it is largest when the current posterior is near the decision threshold.

\section{Distributions}

\subsection{Binomial}

The number of successes in $n$ independent trials with probability $p$ has $\Prob(X = k) = \binom{n}{k}p^k(1-p)^{n-k}$, mean $np$ and variance $np(1-p)$. It is used in the jury and discrimination cases of later chapters.

\subsection{Poisson}

For rare events, the number in a period has $\Prob(X = k) = e^{-\lambda}\lambda^k/k!$, with mean and variance $\lambda$. Poisson derived it in his study of jury verdicts \citep{poisson1837}; see the box ``Poisson and the French juries'' in Chapter~\ref{ch:why}.

\subsection{Beta and gamma}

The beta distribution $\text{Beta}(a,b)$ on $(0,1)$ has density proportional to $\theta^{a-1}(1-\theta)^{b-1}$ and mean $a/(a+b)$. It is conjugate to the binomial: after $k$ successes in $n$ trials the posterior is $\text{Beta}(a+k, b+n-k)$. The gamma distribution with shape $\alpha$ and rate $\beta$ is conjugate to the Poisson: after $k$ events in time $t$ the posterior for the rate is gamma with shape $\alpha + k$ and rate $\beta + t$.

\subsection{Normal}

The normal distribution with mean $\mu$ and variance $\sigma^2$ is conjugate to a normal likelihood with known variance: if $\theta \sim N(m, v)$ and $y\given\theta \sim N(\theta, s^2)$, then $\theta\given y \sim N(m', v')$ with $1/v' = 1/v + 1/s^2$ and $m' = v'(m/v + y/s^2)$. Precisions add and means are precision weighted. This is the formula behind the combination of market and expert evidence in Chapter~\ref{ch:bankruptcy} and the shrinkage in Chapter~\ref{ch:llsv}.

\section{Fallacies}

\begin{description}
\item[Prosecutor's fallacy] Reading $\Prob(E\given\Hd)$ as $\Prob(\Hd\given E)$.
\item[Defence fallacy] Treating evidence as irrelevant because it does not single out the defendant, ignoring its likelihood ratio.
\item[Dershowitz fallacy] Reporting a probability conditioned on less than what is known, such as $\Prob(K\given B)$, the probability that a battered woman is killed by her batterer, when $\Prob(K\given B, M)$, which also conditions on her having been murdered, is required (here $K$ is an event, not background information).
\item[Independence fallacy] Multiplying probabilities of dependent events.
\item[Base rate neglect] Ignoring the prior when interpreting evidence.
\item[Double counting] Adding the weights of items that share a common source.
\item[Hindsight bias] Assessing an ex ante probability using knowledge of the outcome.
\item[Selection neglect] Treating a selected sample, such as tried cases or favourable experts, as representative.
\end{description}

\section{Graphical models}

A Bayesian network factorises a joint distribution as $\prod_i\Prob(X_i\given\mathrm{pa}(X_i))$. In a chain or fork, observing the middle node blocks dependence; in a collider, observing the middle node or a descendant creates it. Two sets of variables are conditionally independent given a third if every path between them is blocked (d-separation).
